% Question 1
\qtag {Objective: 8.A Explain object-oriented programming vocabulary }

\item Which C++ feature allows one interface to support many implementations?
\begin{enumerate}[label=\Alph*.]
\item Polymorphism
\item Recursion
\item Preprocessing
\item Compilation
\end{enumerate}

\begin{solution}
    A
\end{solution}

% Question 2
\qtag {Objective: 6.D Identify languages that are compiled }

\item 
T/F Java source code is compiled before execution. 

\begin{solution}
    T
\end{solution}

% Question 3
\qtag {Objective: 5.A Understand memory allocation }

\item Consider the following code.

\begin{lstlisting}
#include <iostream>
using namespace std;

int main() {
   double cityRainfall = double(2.5);

   cout << "rainfall:" << cityRainfall;

   return 0;
}
\end{lstlisting}

Suppose this program ran multiple times a day.  What best describes this program?  Select all that apply.
\begin{enumerate}[label=\Alph*.]
\item cityRainfall is on the stack.
\item cityRainfall is on the heap.
\item the program has a memory leak.
\end{enumerate}

\begin{solution}
    A
\end{solution}

% Question 4
\qtag {Objective: 7.C Understand linked list insertion }

\item The following function inserts a node into a doubly linked list. Assume that the variable, head, is a pointer to the first node in the list. The member, prev, is a pointer to the previous node, and next is a pointer to the next node. Consider the following pseudocode.

\begin{lstlisting}
function insert(head, newNode):
    if head is null:
        head = newNode
    else:
        current = head
        while current.next is not null:
            current = current.next
        current.next = newNode
        newNode.prev = current
\end{lstlisting}

What is the time complexity of this function?
\begin{enumerate}[label=\Alph*.]
\item O(1)
\item O(n)
\item O(n^2)
\item O(log n)
\item Cannot be determined
\end{enumerate}

\begin{solution}
    B
\end{solution}

% Question 5
\qtag {Objective is a work in progress }

\item Pick the language that is always interpreted, never compiled.
\begin{enumerate}[label=\Alph*.]
\item Python is always interpreted in practice.
\item C++ is compiled.
\item Java uses a virtual machine.
\end{enumerate}

\begin{solution}
    A
\end{solution}
